\textbf{RBX1 is hard.} It is 108 amino acids long. Half of those
residues are intrinsically disordered. The folded RING-H2 core depends
on three structural Zn²⁺ ions held in a cross-brace coordination
{[}@zheng2002scf{]}. In vivo it is the catalytic tip of every cullin-RING
ligase (CRL), recruiting an E2 ubiquitin-conjugating enzyme via its α2
helix and handing ubiquitin onto a substrate organised by the rest of
the complex {[}@petroski2005crl{]}. No designed binder against RBX1
has been published. No reference protein binder exists beyond CUL1
itself. The RING surface is shallow, the tail is disordered, and the
catalytic helix is small.

That is why we picked it. The GEM × Adaptyv 2026 competition follows two
prior Adaptyv community campaigns (EGFR 2024 {[}@adaptyv2024egfr{]},
Nipah-G 2025--26 {[}@adaptyv2026nipah{]}) and escalates the target
difficulty. Submitters could enter any modality (miniproteins, peptides,
single-domain antibodies, scFvs) provided they kept ≤250 aa, ≤75\%
identity to known proteins, and ≤100 sequences per team. The submission
window opened 15 Feb 2026 and closed 26 Mar 2026. Validation ran at the
Adaptyv Foundry through April. Results were announced at the ICLR GEM
workshop on 26 Apr 2026.

This paper reports the wet-lab outcome of the 322-design batch. The
batch mixes methods (RFdiffusion, BoltzGen, Mosaic, BindCraft 2, MoPPIt,
LFM2, hallucination on AF2, and 25 more) and modalities so no single
design philosophy is privileged. The submission table, every ESMFold
prediction, and every SPR replicate are published at ProteinBase under
ODC-ODbL.
